\subsection{Measurement setup}
The data were collected on the ARA wireless living lab~\cite{islam2025design}. A client node is connected to a Starlink user terminal, a OneWeb user terminal, and a 5G network through a Quectel modem. During each session the node runs a ping stream at 5\,Hz to a server, an iperf3 UDP stream that reports throughput, jitter, and loss once per second, and collectors for terminal telemetry (Starlink dish state and satellite geometry, OneWeb SINR, elevation, azimuth and handover events, 5G RSRP, RSRQ and SINR) and local weather. The parsers align all sources on the ping time base, which gives one row every 200\,ms. Figure~\ref{fig:system} shows the setup and where the forecaster sits in the multipath client.

\begin{figure}[t]
  \centering
  \resizebox{0.82\linewidth}{!}{% System model: ARA measurement setup and prediction-aware multipath client.
\begin{tikzpicture}[
  font=\footnotesize, >=Latex,
  sat/.style={draw, fill=blue!12, minimum width=7mm, minimum height=4mm, rounded corners=1pt},
  ut/.style={draw, fill=orange!25, minimum width=14mm, minimum height=6mm, align=center, rounded corners=2pt},
  blk/.style={draw, fill=gray!10, minimum height=7mm, align=center, rounded corners=2pt, text width=20mm},
  pred/.style={draw, fill=green!15, minimum height=7mm, align=center, rounded corners=2pt, text width=22mm},
  lnk/.style={->, thick},
  wl/.style={dashed, thick, blue!60!black}
]
% Satellites
\node[sat] (s1) at (0,4.2) {SL};
\node[sat] (s2) at (1.6,4.6) {SL};
\node[sat] (o1) at (4.0,4.6) {OW};
\node[sat] (o2) at (5.6,4.2) {OW};
\node[font=\scriptsize, align=center] at (0.8,5.15) {Starlink shell\\($\approx$550\,km)};
\node[font=\scriptsize, align=center] at (4.8,5.15) {OneWeb shell\\($\approx$1200\,km)};
\draw[->, gray] (s1.east) -- ++(0.35,0.1);
\draw[->, gray] (o2.east) -- ++(0.35,-0.1);
% UTs and gNB
\node[ut] (sut) at (0.4,2.4) {Starlink\\UT};
\node[ut] (out) at (4.8,2.4) {OneWeb\\UT};
\node[ut, fill=red!15] (gnb) at (2.6,2.4) {5G gNB\\(Quectel UE)};
\draw[wl] (sut) -- (s1); \draw[wl] (sut) -- (s2);
\draw[wl] (out) -- (o1); \draw[wl] (out) -- (o2);
\node[font=\scriptsize, blue!60!black] at (-0.55,3.4) {handover};
% Client node
\node[blk, text width=58mm, minimum height=8mm] (client) at (2.6,0.9) {ARA client node: ping (5\,Hz), iperf3 (1\,Hz), UT gRPC/API telemetry, weather};
\draw[lnk] (sut) -- (sut |- client.north);
\draw[lnk] (gnb) -- (client);
\draw[lnk] (out) -- (out |- client.north);
% Pipeline
\node[blk] (feat) at (-0.6,-0.7) {features $\mathbf{x}_t\in\mathbb{R}^{F}$\\(sampled at 5\,Hz)};
\node[pred] (model) at (2.6,-0.7) {forecaster $f_\theta$\\$\hat{\mathbf{y}}_{t+1:t+H}$};
\node[pred] (cp) at (5.8,-0.7) {conformal\\interval $\pm q_{h,k}$};
\draw[lnk] (client.south -| feat) -- (feat);
\draw[lnk] (feat) -- (model);
\draw[lnk] (model) -- (cp);
\node[blk, text width=58mm, fill=violet!12] (sched) at (2.6,-2.1) {prediction-aware MPQUIC scheduler (liquid tunnel): per-path RTT, jitter, loss forecasts with uncertainty};
\draw[lnk] (cp.south) |- (sched.east);
\draw[lnk] (model.south) -- (sched.north);
\end{tikzpicture}
}
  \caption{System model. Three access links are measured in parallel on the ARA client node. The forecaster maps a window of link-state features to a distribution of future RTT, jitter, and loss, which a conformal step turns into calibrated intervals for the prediction-aware MPQUIC scheduler.}
  \label{fig:system}
\end{figure}

Figure~\ref{fig:sysdetail} gives a more detailed view of the three planes of the system: the network and measurement plane, the data and forecasting plane developed in this report, and the scheduling plane that consumes the forecasts.

\begin{figure}[p]
  \centering
  \rotatebox{90}{\resizebox{0.95\textheight}{!}{% Detailed system model: network and measurement plane, data and forecasting plane, scheduling plane.
\begin{tikzpicture}[
  font=\scriptsize, >={Latex[length=1.6mm]},
  sat/.pic={\fill[gray!60] (-0.32,-0.07) rectangle (-0.12,0.07); \fill[gray!60] (0.12,-0.07) rectangle (0.32,0.07);
            \filldraw[fill=white, draw=black] (-0.11,-0.1) rectangle (0.11,0.1);},
  box/.style={draw, rounded corners=2pt, align=center, inner sep=2.5pt},
  gnd/.style={box, fill=gray!12, minimum height=6mm},
  ue/.style={box, fill=orange!20, minimum width=13mm, minimum height=7mm},
  pstep/.style={box, fill=blue!8, text width=37mm, minimum height=7mm},
  fstep/.style={box, fill=green!12, text width=37mm, minimum height=7mm},
  sch/.style={box, fill=violet!12, text width=30mm, minimum height=7mm},
  panel/.style={draw=gray!70, dashed, rounded corners=4pt, inner sep=5pt},
  ptitle/.style={font=\footnotesize\bfseries},
  ul/.style={dashed, thick, blue!60!black},
  fl/.style={densely dotted, very thick, blue!60!black},
  wl/.style={thick, ->}
]
% ===================== (a) network and measurement plane =====================
\node[ptitle, anchor=west] (ta) at (-0.5,8.75) {(a) Network and measurement plane};
\draw[gray!60, thick] (-0.3,6.3) .. controls (2.5,6.75) and (4.5,6.75) .. (7.6,6.3);
\draw[gray!60, thick] (-0.3,7.5) .. controls (2.5,8.0) and (4.5,8.0) .. (7.6,7.5);
\node[gray!60!black, anchor=west, align=left] at (6.0,6.15) {Starlink\\$\approx$550\,km};
\node[gray!60!black, anchor=west, align=left] at (6.0,7.35) {OneWeb\\$\approx$1200\,km};
\pic at (0.7,6.5) {sat}; \pic at (2.1,6.63) {sat}; \pic at (3.5,6.66) {sat};
\pic at (3.3,7.82) {sat}; \pic at (5.0,7.82) {sat};
\node at (0.7,6.83) {S$_1$}; \node at (2.1,6.96) {S$_2$}; \node at (3.5,6.99) {S$_3$};
\node at (3.3,8.15) {O$_1$}; \node at (5.0,8.15) {O$_2$};
\draw[<->, red!70!black] (1.05,6.56) -- (1.75,6.61);
\node[red!70!black, font=\tiny, align=center] at (1.4,7.2) {handover\\(15\,s reallocation)};
\draw[<->, red!70!black] (3.65,7.86) -- (4.65,7.88) node[midway, above, font=\tiny] {handover};
% ground segment
\node[gnd] (gs) at (1.7,4.4) {Starlink\\GS + PoP};
\node[gnd] (gw) at (4.6,4.4) {OneWeb\\gateway};
\node[gnd] (core) at (5.25,3.55) {5G core};
\node[box, fill=gray!25, minimum width=14mm] (inet) at (7.0,5.0) {Internet};
\node[box, fill=yellow!20] (srv) at (7.0,4.0) {server:\\ping echo,\\iperf3};
\draw[fl] (gs) -- (2.1,6.5);
\draw[fl] (gw) -- (5.0,7.7);
\draw (gs) -- (inet); \draw (gw) -- (inet); \draw (core) -- (inet); \draw (inet) -- (srv);
% user side
\node[ue] (sut) at (0.45,1.75) {Starlink\\UT};
\node[ue] (out) at (3.05,1.75) {OneWeb\\UT};
\node[ue] (qe) at (5.95,1.75) {Quectel\\5G UE};
\node[gnd] (gnb) at (5.95,2.7) {gNB};
\draw[ul] (sut) -- (0.7,6.38);
\draw[ul] (out) -- (3.3,7.7);
\draw[ul] (qe) -- (gnb); \draw (gnb) -- (core);
\node[box, fill=orange!8, minimum width=70mm, minimum height=8mm] (cn) at (3.1,0.55) {ARA client node: MPQUIC client, collectors, forecaster};
\draw[thick] (sut) -- (sut |- cn.north); \draw[thick] (out) -- (out |- cn.north); \draw[thick] (qe) -- (qe |- cn.north);
\draw[ul] (1.05,3.15) -- (1.5,3.15) node[right] {user link};
\draw[fl] (1.05,2.8) -- (1.5,2.8) node[right] {feeder link};
\node[panel, fit={(ta)(cn)(srv)(sut)(7.9,8.2)}] {};

% ===================== (b) data and forecasting plane =====================
\node[ptitle, anchor=west] (tb) at (8.45,8.75) {(b) Data and forecasting plane};
\node[pstep] (col) at (10.6,7.75) {\textbf{Collectors per link $p$}: ping RTT (5\,Hz), iperf3 jitter, loss, rate (1\,Hz), terminal telemetry, weather};
\node[pstep] (ali) at (10.6,6.6) {\textbf{Causal alignment} on a 200\,ms grid: interpolated pings $\rightarrow$ last real ping (audit item A3)};
\node[pstep] (feat) at (10.6,5.45) {\textbf{Features} $\mathbf{x}_t\in\mathbb{R}^F$ + causal context statistics, $z$-scored, clipped};
\begin{scope}[shift={(8.7,3.45)}]
  \draw[->] (0,0) -- (3.9,0);
  \foreach \x in {0.1,0.25,...,2.0} \fill[black] (\x,{0.35+0.18*sin(220*\x)+0.1*(\x>1.3)}) circle (0.6pt);
  \draw[decorate, decoration={brace, amplitude=3pt, mirror}] (0.05,-0.08) -- (2.0,-0.08) node[midway, below=2pt, font=\tiny] {$\mathbf{X}_t$: $L=30$ rows (6\,s)};
  \draw[dashed, gray] (2.05,-0.1) -- (2.05,1.05) node[above, font=\tiny, black] {$t$};
  \fill[green!30, opacity=0.7] (2.1,0.48) -- (3.8,0.2) -- (3.8,1.15) -- (2.1,0.62) -- cycle;
  \fill[green!55, opacity=0.7] (2.1,0.52) -- (3.8,0.42) -- (3.8,0.85) -- (2.1,0.58) -- cycle;
  \draw[green!40!black, thick] (2.1,0.55) -- (3.8,0.6);
  \draw[decorate, decoration={brace, amplitude=3pt, mirror}] (2.1,-0.08) -- (3.8,-0.08) node[midway, below=2pt, font=\tiny] {$h=1..15$ (0.2--3\,s)};
\end{scope}
\node[fstep] (enc) at (10.6,2.45) {\textbf{DS2S}: GRU encoder $\rightarrow$ attention $\rightarrow$ decoder anchored on $\mathbf{y}_t$};
\node[fstep] (head) at (10.6,1.3) {\textbf{Heads}: $\hat\mu_{h,k}$, $\hat q_{\tau,h,k}$, $\tau\in\{.025,.25,.5,.75,.975\}$, loss MSE + pinball};
\node[fstep] (cqr) at (10.6,0.15) {\textbf{CQR}: $[\hat q_{.025}-c_{h,k},\,\hat q_{.975}+c_{h,k}]$, $c_{h,k}$ from calibration sessions};
\draw[wl] (col) -- (ali); \draw[wl] (ali) -- (feat); \draw[wl] (feat) -- (10.6,4.72); \draw[wl] (10.6,3.05) -- (enc);
\draw[wl] (enc) -- (head); \draw[wl] (head) -- (cqr);
\node[panel, fit={(tb)(col)(cqr)(8.55,3.0)(12.65,3.0)}] {};
\draw[wl, orange!70!black] (cn.east) to[out=0, in=180] node[pos=0.12, below right, font=\tiny, black] {measurements} (col.west);

% ===================== (c) scheduling plane =====================
\node[ptitle, anchor=west] (tc) at (13.55,8.75) {(c) Scheduling plane};
\node[sch] (fc) at (15.45,7.5) {per path $p$, step $h$: median $\hat q_{.5}$, mean $\hat\mu$, bound $u_{p,h}=\hat q_{.975}+c$};
\node[sch] (app) at (15.45,6.1) {real-time application packets with deadline $D$};
\node[sch] (dec) at (15.45,4.6) {\textbf{Prediction-aware MPQUIC scheduler}: send on $p$ with $u_{p,h}\le D$, else duplicate on two paths};
\node[box, fill=orange!20, minimum width=8mm] (p1) at (14.25,2.9) {SL};
\node[box, fill=orange!20, minimum width=8mm] (p2) at (15.45,2.9) {OW};
\node[box, fill=orange!20, minimum width=8mm] (p3) at (16.65,2.9) {5G};
\draw[wl] (fc) -- (app); \draw[wl] (app) -- (dec);
\draw[wl] (dec) -- (p1); \draw[wl] (dec) -- (p2); \draw[wl] (dec) -- (p3);
\node[align=center, font=\tiny] (sf) at (15.45,2.2) {QUIC subflows on the three\\access links of (a)};
\node[sch, fill=gray!10] (fb) at (15.45,0.95) {ACKs and new measurements update $\mathbf{X}_{t+1}$ every 200\,ms};
\draw[wl, dashed] (sf) -- (fb);
\node[panel, fit={(tc)(fc)(fb)(p1)(p3)(cqr.south -| fb.south)}] {};
\draw[wl, green!40!black] (cqr.east) to[out=0, in=180] node[pos=0.1, below right, font=\tiny, black, align=left] {forecasts +\\intervals} (fc.west);
\end{tikzpicture}
}}
  \caption{Detailed system model. (a) Network and measurement plane: Starlink (about 550\,km, user reallocation every 15\,s), OneWeb (about 1200\,km, bent pipe), and 5G paths from the ARA client node to the measurement server. (b) Data and forecasting plane: causal alignment, feature window, distributional forecaster, and CQR calibration. (c) Scheduling plane: per-path forecasts and calibrated bounds drive the prediction-aware MPQUIC scheduler, and new measurements update the window every 200\,ms.}
  \label{fig:sysdetail}
\end{figure}

\subsection{Forecasting problem}
For a link, let $\mathbf{x}_t\in\mathbb{R}^{F}$ be the feature vector at row $t$ and $\mathbf{y}_t\in\mathbb{R}^{K}$ the $K=3$ targets (RTT in ms, iperf jitter in ms, iperf loss in \%). Given the window $\mathbf{X}_t=(\mathbf{x}_{t-L+1},\dots,\mathbf{x}_t)$ with $L=30$ rows, the model predicts $\mathbf{y}_{t+h}$ for $h=1,\dots,H$ with $H=15$. Because rows are 200\,ms apart, the input covers 6\,s of history and the horizon spans 0.2\,s to 3\,s. The Year 2 report labeled the horizon steps as seconds (\texttt{t+1s} to \texttt{t+15s}). They are steps of 200\,ms, which matters when the forecasts are matched to scheduler time constants.

\subsection{Datasets}
Table~\ref{tab:data} lists the six datasets. Sessions (one CSV file each) are split 70/15/15 into training, validation, and test sets by a random permutation with seed 42, exactly as in Year 2. All windows from one session therefore stay in one split, and the test sessions are never seen during training, model selection, or calibration.

\begin{table}[t]
\centering
\caption{Datasets. Windows are 30 input rows plus 15 target rows at 5\,Hz. For 5G, every fourth training window is used (Section~\ref{sec:protocol}). Features: Year 2 set / Year 3 set.}
\label{tab:data}
\small
\begin{tabular}{lrrrrrr}
\toprule
Link & Sessions (tr/va/te) & Train windows & Val windows & Test windows & Features \\
\midrule
OneWeb DL   & 45 / 9 / 11   & 117{,}558 & 22{,}655  & 27{,}492  & 9 / 31 \\
OneWeb UL   & 45 / 9 / 11   & 115{,}847 & 22{,}751  & 28{,}125  & 9 / 31 \\
Starlink DL & 37 / 7 / 9    & 105{,}906 & 20{,}025  & 25{,}838  & 19 / 41 \\
Starlink UL & 37 / 7 / 9    & 106{,}165 & 20{,}090  & 25{,}666  & 19 / 41 \\
5G DL       & 266 / 57 / 58 & 194{,}725$^\dagger$ & 166{,}266 & 169{,}204 & 10 / 32 \\
5G UL       & 265 / 56 / 58 & 193{,}502$^\dagger$ & 162{,}894 & 168{,}519 & 10 / 32 \\
\bottomrule
\end{tabular}

\smallskip
\footnotesize $^\dagger$ after keeping every fourth window (778{,}898 and 774{,}006 before subsampling).
\end{table}
